\section{Datasheet}\label{app:datasheet}
Following \citet{gebru2021}.

\paragraph{Motivation.} The benchmark measures how well the maximal growth rate of a microbe can be predicted from its
genome when the closest labeled relative is at a known taxonomic distance, which is the situation of genomes assembled
from environmental DNA. It was created by the authors of this paper; no external funding is reported.

\paragraph{Composition.} \NGenomes genomes of \NSpecies species (\NArchaea archaeal), each with a maximal doubling time in
hours at the optimal growth temperature, the growth temperature, and the doubling time standardized to \TRefC\,\textdegree C
(Section~\ref{sec:data}); \NSlow have a doubling time above one day. Inputs are computed from the genome sequence alone.
Each genome carries its GTDB R\GTDBRelease taxonomy and species group, which define the splits. \NExcluded labels were
removed with written evidence and are released as a separate file, so that their effect can be studied. The data contain
no information about people.

\paragraph{Collection.} Labels come from the growth-rate table distributed with gRodon \citep{weissman2021}, derived from
the compilation of \citet{madin2020}, and from \NLiteratureWord strains whose doubling time was read in the primary
paper. Genomes were downloaded from NCBI by accession. Most labels are species values, not values for the sequenced strain.

\paragraph{Preprocessing.} Names were mapped to GTDB; rows were joined into species groups; genes were predicted with
pyrodigal; inputs were computed as in Section~\ref{sec:data}; labels were audited with Algorithm~\ref{alg:audit}. The raw
genome sequences are not redistributed; the scripts download them.

\paragraph{Uses.} Benchmarking genome-to-phenotype prediction under taxonomic distribution shift, and uncertainty
estimation under that shift. It should not be used to infer growth rates in the environment, which are usually much lower
than maximal rates in culture, nor as a census of slow growers.

\paragraph{Distribution and maintenance.} Code and curated labels are in a public repository
(link withheld for anonymous review); the derived tables, result files and predictions form an archive
built by \texttt{paper/scripts/make\_archive.py} for deposit at Zenodo. Code is released under Apache-2.0 and the derived
data under CC BY 4.0; the upstream sources keep their own terms. Errors can be reported as repository issues, and label
corrections are added to the exclusion file with their evidence.
